% Appendix A: rank-full configurations. Source of truth: v3, Appendix A.
\section{Rank-full configurations}\label{app:rankfull}
% v3: Appendix A

This appendix proves the moment bound for rank-full clusters of generalized integers in a short window
(Theorem~\ref{thm:R}). It is used in the proof of the cluster event (Proposition~\ref{prop:ecl}). Its machinery
(Lemmas~\ref{lem:volume}--\ref{lem:SC} and Theorem~\ref{thm:svl}) is also used in the Cover Bound
(Proposition~\ref{prop:cover}). Probabilistically, the argument rests only on the standing facts
\ref{eq:sF1}--\ref{eq:sF4} of Section~\ref{sec:queue-standing}:
\begin{itemize}
\item the slot count \ref{eq:sF1};
\item the jitter density \ref{eq:sF2};
\item the harmonic sums \ref{eq:sF3p};
\item the sure Mertens bound \ref{eq:sF4}.
\end{itemize}
These are carried by progressive good events.

The idea is to integrate out the jitter of $k$ suitably chosen primes (the \emph{pivots}), one for each element of the
cluster. Each pivot is integrated against its conditional density. Every other prime (a \emph{free} prime) is summed over
its realised values by a sure bound.

Two difficulties have to be overcome:
\begin{itemize}
\item a prime that occurs in several elements constrains several of them at once;
\item a free prime may be decided after some pivot, so its realisation is not measurable at that pivot's decision time.
\end{itemize}
The \emph{Volume Lemma} (Lemma~\ref{lem:volume}) handles the first: integrating all pivots jointly by a change of
variables replaces per-constraint widths by an exact Jacobian. The \emph{Pivot-Integration Lemma} (Lemma~\ref{lem:pil})
and the \emph{Staged Volume Lemma} (Theorem~\ref{thm:svl}) handle the second: integrate the pivots in stages, and bound
the free primes decided between stages by sure sums, each against a weight that is harmonic in some row.

\subsection{Setting}\label{sec:A-setting}
% v3: Appendix A, A.0

Throughout, $L=\log x$, $\ell=\log\log x$ and $\Lambda=\log X_0\ge 10^{8}$, as in Section~\ref{sec:queue-standing}.
Fix a dyadic height $z\ge X_0/2$ and a window $J\subset[z,2z]$ of length $\delta\le 1$, and put $\eps:=\delta/z$.

\begin{definition}[Configuration types]\label{def:config-type}
% v3: Appendix A, A.0
A \emph{configuration type} $\tau$ fixes:
\begin{enumerate}[label=(\alph*)]
\item $k$ rows $n_1,\dots,n_k\in J$, which are elements of $\NN$;
\item the jittered primes occurring in them (the \emph{columns}), listed in decision order;
\item the integer exponent matrix $E=(E_{ij})$, where $E_{ij}$ is the exponent of the $j$-th column in $n_i$;
\item for each row, its $X_0$-smooth tail $d_i$, i.e.\ the product of its rational primes $<X_0$.
\end{enumerate}
The type $\tau$ is \emph{rank-full} if $\operatorname{rank}E=k$. For a column $p$, write $e_p\in\Z_{\ge0}^{k}$ for its
column of $E$.
\end{definition}

\paragraph{Greedy pivots.}
Run through the columns latest-first, and declare $p$ a \emph{pivot} if and only if
\[
  e_p\notin\operatorname{span}\{e_q:\ q\text{ a pivot decided later than }p\}.
\]
Let $P$ be the set of pivots. For a rank-full type, $|P|=k=\operatorname{rank}E$ and the square matrix $E_P$ (the
columns of $E$ indexed by $P$) is invertible. The remaining columns are \emph{free}. For $i\le k$ write
\[
  n_i=c_i\prod_{j\in P}q_j^{E_{ij}},
\]
where $q_j$ is the value of the $j$-th pivot and $c_i\in\NN$ collects the free jittered primes of $n_i$ together with its
tail $d_i$.

\paragraph{Runs and blocks.}
Split $P$ into maximal runs $P_1,\dots,P_T$ of consecutive pivots in decision order, latest first. For $t<T$, the
\emph{late free block} $F_t$ consists of the free columns decided between $P_t$ and $P_{t+1}$. The \emph{early block}
$F_T$ consists of the free columns decided before all pivots. Write $P_{\le t}=P_1\cup\dots\cup P_t$, and similarly
$F_{\le t}$. By construction of the greedy pivots,
\begin{equation}\label{eq:A-span}
  e_f\in\operatorname{span}\{e_p:\ p\in P_{\le t}\}\qquad\text{for every }f\in F_{\le t}.
\end{equation}
Indeed, a free column $f$ decided after $P_{t+1}$ lies in the span of the pivots decided later than $f$, and these
belong to $P_{\le t}$.

\paragraph{Slot positions.}
Pivots are distinct primes, one per slot. Slot spacing is $\ge 0.99$, and the value $q$ of a jittered prime differs
from its slot $s_q$ by at most $H$, the jitter width at the slots in question (Section~\ref{sec:queue-standing}).
Hence, whatever the labels are, the $j$-th earliest pivot has value at least
\begin{equation}\label{eq:A-tau}
  \tau_j:=\max\bigl(X_0,\ X_0+0.99(j-1)-2H\bigr).
\end{equation}
Here and below, ``a jittered prime is $\ge X_0$'' is read with $X_0^{-}=X_0-H(X_0)/2$ in place of $X_0$, as in
Section~\ref{sec:queue-standing}. Every such use has a margin far larger than $H(X_0)/X_0$.

\paragraph{Probabilistic set-up.}
$\mathcal F_{s-}$ is the information available before the decision at slot~$s$. $\Good(<s)$ is the good event of
Definition~\ref{def:good}, and $G_{<s}$ abbreviates it. For a window $J$, $t_J$ denotes the time at which the
cluster event of $J$ is examined; it satisfies $t_J\le t(E1)$ (Definition~\ref{def:good}). Every bound below is proved
on $\Good(<t_J)$, which is carried progressively along the pivots. No queue bound at heights inside the window is
used quantitatively.

\subsection{The Volume Lemma and the pattern sum}\label{sec:A-volume}

\begin{lemma}[Volume Lemma]\label{lem:volume}
% v3: Lemma A.1
Fix the cofactors $c_1,\dots,c_k$. The map $q_P=(q_j)_{j\in P}\mapsto(n_i)_{i\le k}$, $n_i=c_i\prod_jq_j^{E_{ij}}$,
is injective on $\R_{+}^{k}$, and
\begin{equation}\label{eq:volume}
  \operatorname{vol}\bigl\{q_P\in\R_+^k:\ n_i\in J\ \forall i\bigr\}
  =\int_{J^k}\frac{\prod_jq_j}{|\det E_P|\,\prod_in_i}\,dn,
  \qquad
  \prod_jq_j=\frac{\prod_i(n_i/c_i)}{\prod_jq_j^{\,s_j-1}},
\end{equation}
where $s_j:=\sum_iE_{ij}\ge1$ is the $j$-th column sum. In particular, if all $q_j\ge X_0$, then the volume is at most
$\delta^k/\bigl(|\det E_P|\prod_ic_i\,X_0^{M-k}\bigr)$, with $M:=\sum_{i,j}E_{ij}$.
\end{lemma}

\begin{proof}
Pass to logarithmic coordinates $x_j=\log q_j$. Then $\log n_i=\log c_i+\sum_jE_{ij}x_j$, an affine map with
invertible linear part $E_P$, so the map is injective. Its Jacobian matrix is $\partial n_i/\partial q_j=E_{ij}n_i/q_j$,
whose determinant is $\det(E_P)\prod_in_i/\prod_jq_j$. The change of variables $q_P\mapsto n$ therefore gives the
integral in \eqref{eq:volume}.

For the second identity, multiply the relations $n_i/c_i=\prod_jq_j^{E_{ij}}$ over $i$. This gives
$\prod_jq_j^{s_j}=\prod_i(n_i/c_i)$, hence $\prod_jq_j=\prod_i(n_i/c_i)\big/\prod_jq_j^{s_j-1}$. No column of the
invertible matrix $E_P$ vanishes, so $s_j\ge1$. If all $q_j\ge X_0$, then $\prod_jq_j^{s_j-1}\ge X_0^{M-k}$. Hence the
integrand is at most $1/\bigl(|\det E_P|\prod_ic_iX_0^{M-k}\bigr)$, and $J^k$ has volume $\delta^k$.
\end{proof}

\begin{remark}
Lemma~\ref{lem:volume} involves no per-constraint widths. Bounding the constraints one at a time by Gaussian elimination
would combine them into integer combinations whose $\ell^1$-norms can grow exponentially in $k$. The Jacobian avoids this
entirely.
\end{remark}

\begin{lemma}[Pattern sum]\label{lem:pattern}\label{lem:patternsum}
% v3: Lemma A.2
Summing over the invertible nonnegative integer matrices $E_P$ with columns in decision order,
\[
  \frac1{k!}\sum_{E_P}\prod_j\tau_j^{-(s_j-1)}\ \le\ \prod_j\bigl(1-1/\tau_j\bigr)^{-k}
  \ \le\ e^{(3H(X_0)+2)k/X_0}\bigl(1+k/X_0\bigr)^{1.02k}.
\]
Since $H(X_0)/X_0\le e^{-10^{7}}$, the right side is at most $e^{k\,e^{-10^{6}}}\bigl(1+k/X_0\bigr)^{1.02k}$.
% CORRECTIONS-vs-v3.md entry 1: v3 states e^{2k/X0}; corrected constant e^{(3H(X0)+2)k/X0}.
\end{lemma}

\begin{proof}
By the Leibniz expansion, an invertible $E_P$ has a nonzero term $\prod_iE_{i,\sigma(i)}$ for some permutation $\sigma$.
So its support contains the permutation matrix of $\sigma$. Subtract that permutation matrix. The rest is a nonnegative
integer matrix $X$, and the column sums satisfy $s_j=1+\sum_iX_{ij}$. Choosing one such $\sigma$ for each $E_P$ makes the
map $E_P\mapsto(\sigma,X)$ injective. Hence
\[
  \sum_{E_P}\prod_j\tau_j^{-(s_j-1)}\le k!\prod_j\sum_{m\ge0}\binom{k+m-1}{m}\tau_j^{-m}=k!\prod_j(1-1/\tau_j)^{-k},
\]
since a column of extra mass $m$ can be distributed over the $k$ rows in $\binom{k+m-1}{m}$ ways. Next,
$(1-1/\tau)^{-k}\le\exp\bigl(k/(\tau-1)\bigr)$. The indices with $\tau_j=X_0$ number at most $2H/0.99+1$, and the others
satisfy $\tau_j\ge X_0+0.99i$ for distinct $i\ge1$. Hence
\[
  \sum_{j\le k}\frac1{\tau_j-1}\le\frac{3H}{X_0}+\frac1{X_0-1}+\sum_{i=1}^{k}\frac1{X_0-1+0.99i}
  \le\frac{3H}{X_0}+\frac1{X_0-1}+1.02\log\Bigl(1+\frac{k}{X_0}\Bigr).
\]
Since $1/(X_0-1)\le2/X_0$, exponentiating gives
$\prod_j(1-1/\tau_j)^{-k}\le e^{(3H(X_0)+2)k/X_0}(1+k/X_0)^{1.02k}$. Here $H=H(X_0)$, the jitter width at the slots
near $X_0$, where the pivots with $\tau_j=X_0$ lie.
\end{proof}

\subsection{The Pivot-Integration Lemma}\label{sec:A-pil}

Let $X_\tau$ be the number of realised configurations of type $\tau$ satisfying all window and feasibility
constraints. These are $k$-tuples of distinct elements of $\NN\cap J$ with the prescribed exponent pattern. Let
$p_1,p_2,\dots$ be the pivots, latest first.

\begin{lemma}[Pivot-Integration Lemma]\label{lem:pil}\label{lem:PIL}
% v3: Lemma A.3
Define $S_\tau$ by the following latest-first recursion, starting from $\Psi:=\prod_i\mathbf 1(n_i\in J)$:
\begin{enumerate}[label=(\roman*)]
\item at the pivot $p_t$: $\Psi\leftarrow\int\Psi\,dq_{p_t}$ (Lebesgue measure; the window restriction of $p_t$ is
  dropped);
\item at a late block $F_t$: $\Psi\leftarrow$ any bound for the realised sum $\sum_{F_t}\Psi$ that holds on
  $G_{<s_{p_t}}$ and is a function of variables decided before the next pivot $p_{t+1}$. Examples are the
  \emph{sup-sum} $\sup_{\omega\in G_{<s_{p_t}}}\sum_{F_t}\Psi$, or a carried Euler factor
  (Lemma~\ref{lem:fac}\,(S3));
\item at the early block and the tails: ordinary realised sums, taken at the end.
\end{enumerate}
Then
\[
  \E\bigl[X_\tau\,\mathbf 1(\Good(<t_J))\bigr]\le\theta^{k}\,\sup S_\tau,\qquad\theta=1+1/L .
\]
\end{lemma}

\begin{proof}
Write $X_\tau$ as a sum, over the slot labels of the pivots and over the realised free data, of
$\prod_i\mathbf 1(n_i\in J)$. Every pivot is decided before $t_J$. Along the pivots, latest first, the good events are
nested:
\[
  \mathbf 1(G_{<t_J})\le\mathbf 1(G_{<s_{p_1}})\le\mathbf 1(G_{<s_{p_2}})\le\cdots .
\]
Each $\mathbf 1(G_{<s_{p}})$ is $\mathcal F_{s_p-}$-measurable (Definition~\ref{def:good}). We induct along the decision
order, latest first, keeping the invariant that after each step the pending integrand is a nonnegative function of the
variables not yet processed.

\emph{Pivot step.} Let $p$ be the current pivot, with slot $s$. Everything still pending is a function of $q_p$ and of
variables decided before $s$. By \ref{eq:sF2}, the position of the prime decided at slot $s$ has conditional density
$\le1/H(s)$ on its window $W_s$, given $\mathcal F_{s-}$. Hence
\[
  \E\bigl[\Psi\,\mathbf 1(G_{<s})\bigm|\mathcal F_{s-}\bigr]\le\mathbf 1(G_{<s})\,\frac1{H(s)}\int_{W_s}\Psi\,dq_p .
\]
The existence of a prime at slot $s$ is bounded by $1$ and dropped. Summing over the possible slots $s$ of $p$,
\[
  \sum_s\frac1{H(s)}\int_{W_s}\Psi\,dq=\int\Psi(q)\sum_s\frac{\mathbf 1(q\in W_s)}{H(s)}\,dq\le\theta\int\Psi\,dq .
\]
This uses \ref{eq:sF1}: at most $H(t)(1+\kappa_+/\log t)+3$ windows contain a given point $t$, and $H$ varies
negligibly over a window, so the inner sum is $\le1+\kappa_+/L+3/H\le\theta$.

\emph{Block step.} At a late block $F_t$, the pending $\Psi$ is a nonnegative function of the realised block data and of
earlier variables. Replacing the realised sum by a bound that holds on $G_{<s_{p_t}}$, and that is a function of
variables decided before $p_{t+1}$, keeps the inequality on the event carried at this point. The result is
$\mathcal F_{s_{p_{t+1}}-}$-measurable, which is exactly what the next pivot step needs.

\emph{Order of decision versus order of value.} Slots order the decisions, while the prime values carry jitter.
Primes in $(q-H,q]$ that are decided after the slot of $q$ number at most $H+3$, each of value $\ge X_0$. Wherever a sure
bound is indexed by value rather than by slot, this costs a factor $\le e^{(H+3)/X_0}$ at each end of a chain. It is
absorbed in $\theta$ and in Lemma~\ref{lem:fac}.

After the last pivot, the early block and the tails are summed over their realised values, which are
$\mathcal F$-measurable before every pivot. Taking expectations completes the induction.
\end{proof}

\subsection{Factorisation and sure sums}\label{sec:A-fac}

\begin{lemma}[Factorisation; sure sums]\label{lem:fac}\label{lem:FAC}
% v3: Lemma A.4
\begin{enumerate}[label=(\arabic*),leftmargin=*]
\item[\textup{(FAC)}] Let $f$ be a free column and $P_{\mathrm{late}}$ the set of pivots decided later than $f$, so
  $e_f=E_{P_{\mathrm{late}}}v_f$ for a unique $v_f$ by \eqref{eq:A-span}. After the pivots in $P_{\mathrm{late}}$ have
  been integrated,
  \[
    \Psi(q_f,\text{earlier})\le q_f^{-\beta_f}\,\mathbf 1\bigl(q_f\in C_\eps(\text{earlier})\bigr)\,\Psi_0(\text{earlier}),
    \qquad\beta_f:=\mathbf 1^{\mathsf T}v_f .
  \]
  Here $C_\eps$ is the projection of the feasibility caps (later pivots $\ge\tau_j$, and decision order) onto the
  $q_f$-axis, over the later-pivot values compatible with the window constraints; this widens the caps by relative
  $O(\eps)$. The thin window constraints are invariant under the shift $x_f\mapsto x_f+s$,
  $x_{P_{\mathrm{late}}}\mapsto x_{P_{\mathrm{late}}}-s\,v_f$ (with $x=\log q$). So no late free prime ever enters a
  thin indicator.
\item[\textup{(S1$'$)}] Let $c$ range over cofactors all of whose primes are $\le P:=s_{p_t}+H$. Then for $\beta\ge1$,
  \[
    \sum_c c^{-\beta}\le\mathrm{Eu}(P),\qquad \mathrm{Eu}(x):=\prod_{\text{realised }p\le x}(1-1/p)^{-1}.
  \]
\item[\textup{(S2$'$)}] If $\beta<1$ and $c\le Y$, then $\sum_cc^{-\beta}\le Y^{1-\beta}\,\mathrm{Eu}(P)$. In particular
  $\#\{c\le Y:\ P^+(c)\le P\}\le Y\,\mathrm{Eu}(P)$.
\item[\textup{(S3)}] \emph{(Euler chains, indexed by slot.)} Put
  $\mathrm{Eu}_s(s):=\prod_{\text{realised }p\text{ of slot}\le s}(1-1/p)^{-1}$. A cofactor with harmonic weight that is
  split by the pivots into slot ranges telescopes to $\mathrm{Eu}_s(\mathrm{top})/\mathrm{Eu}_s(\mathrm{bottom})$. The
  factor $1/\mathrm{Eu}_s(s_p)$ may be carried as a pending factor: it is $\mathcal F_{s_p-}$-measurable. Fuzz between
  slot and value appears only at the two ends of the chain. On $\Good$, the sure Mertens bound \ref{eq:sF4}
  (Lemma~\ref{lem:cm}\,(iv), applied once per chain as in Remark~\ref{rem:cm-chain}) bounds the whole chain by
  \[
    \frac{\log a_0}{\log a_m}\,e^{2.20/\log a_m+1/X_0},
  \]
  where $a_m<a_0$ are its end values.
\end{enumerate}
\end{lemma}

\begin{proof}
\emph{(FAC).} The window constraints read $\log n_i=\log c_i+\sum_{p}E_{ip}x_p\in\log J$. Since
$e_f=E_{P_{\mathrm{late}}}v_f$, the shift $x_f\mapsto x_f+s$, $x_{P_{\mathrm{late}}}\mapsto x_{P_{\mathrm{late}}}-sv_f$
leaves every $\sum_pE_{ip}x_p$ unchanged. In the integral over the later pivots, substitute
$y:=x_{P_{\mathrm{late}}}+v_fx_f$. The window indicators then no longer involve $x_f$. The Jacobian
$\prod_{p\in P_{\mathrm{late}}}q_p=e^{\mathbf 1^{\mathsf T}x_{P_{\mathrm{late}}}}=e^{\mathbf 1^{\mathsf T}y}\,q_f^{-\beta_f}$
produces the factor $q_f^{-\beta_f}$. The remaining dependence on $x_f$ is through the feasibility caps,
$y_p-v_{f,p}x_f\ge\log\tau_p$ and the decision order. These are half-spaces in $x_f$ for each $y$. Bounding them by their
projection over the $y$-values compatible with the window constraints gives $\mathbf 1(q_f\in C_\eps)$.

\emph{(S1$'$)--(S2$'$).} For $\beta\ge1$, $\sum_cc^{-\beta}\le\sum_cc^{-1}\le\prod_{p\le P}(1-1/p)^{-1}=\mathrm{Eu}(P)$,
since every $c$ is a product of realised primes $\le P$ (rational primes $<X_0$ included). For $\beta<1$ and $c\le Y$,
$c^{-\beta}=c^{-1}c^{1-\beta}\le Y^{1-\beta}c^{-1}$, and the same Euler product bounds $\sum c^{-1}$.

\emph{(S3).} The sum of $1/c$ over cofactors whose primes have slots in a range $(a,b]$ is at most
$\prod_{\text{slot}\in(a,b]}(1-1/p)^{-1}=\mathrm{Eu}_s(b)/\mathrm{Eu}_s(a)$. Consecutive ranges are disjoint, so the
product over the pieces of a chain telescopes exactly to $\mathrm{Eu}_s(\mathrm{top})/\mathrm{Eu}_s(\mathrm{bottom})$.
At a pivot $p$ the carried factor $1/\mathrm{Eu}_s(s_p)$ involves only primes of slot $<s_p$, so it is
$\mathcal F_{s_p-}$-measurable. The next piece contributes $\mathrm{Eu}_s(s_p)$, which cancels it.

At the two ends, compare slots with values: $\mathrm{Eu}_s(\mathrm{top})/\mathrm{Eu}_s(\mathrm{bottom})\le
\mathrm{Eu}(\mathrm{top}+H)/\mathrm{Eu}(\mathrm{bottom}-H)$. By the sure local Mertens bound of Lemma~\ref{lem:cm}\,(iv) (once per chain; Remark~\ref{rem:cm-chain}),
$\sum_{y<p\le x}1/p\le\log(\log x/\log y)+2.20/\log y$ on $\Good$. Moreover
$\sum_{p\ge X_0}\bigl(-\log(1-1/p)-1/p\bigr)\le\sum_{p\ge X_0}p^{-2}\le1/X_0$. The end fuzz $H$ is absorbed as in
Lemma~\ref{lem:pil}.
\end{proof}

\begin{remark}[Widened caps; an optional route]\label{rem:widened}
% v3: Lemma A.4, Remark (widened caps; optional)
Widen the caps of (FAC) by the pinned range $\rho_t$ of the later pivots, instead of by $O(\eps)$. Then they become
functions of outer variables only, the sup-sums commute with the later integrals, and the recursion of
Lemma~\ref{lem:pil} becomes a monotone functional of the final summand. The widened box costs a factor $e^{O(k\rho_t)}$.
Since $\|E_P^{-1}\|_\infty\le k\,m^{k-1}$, which is sharp up to constants, with $m:=\log2z/\Lambda$ the maximal exponent
sum of a row, this is harmless when
\[
  (k-1)\log m+O(\log k)\ \ll\ \log(z/\delta).
\]
The main line does not use this route. It uses the stagewise form (Theorem~\ref{thm:svl} with Lemma~\ref{lem:SC}),
which needs no conditioning hypothesis.
\end{remark}

\subsection{Nested stage rows}\label{sec:A-SC}

Let $B'$ be a rank-full set of $r$ rows, with exponent matrix $E$ on its jittered primes, and with pivots, runs and blocks
as in Section~\ref{sec:A-setting}. A \emph{cover} is a set $C\subseteq B'$ of rows such that every prime occurring in
$B'$ occurs in some row of $C$. For each $t$, let $M_t$ be the row matroid of $E_{\cdot,P_{\le t}}$, i.e.\ the matroid on
the rows $y\in B'$ in which a set of rows is independent iff their restrictions to the columns $P_{\le t}$ are linearly
independent.

\begin{lemma}[Nested stage rows]\label{lem:SC}\label{lem:sc}
% v3: Lemma A.7 (Lemma SC)
Order the rows so that those of $C$ come before those of $B'\setminus C$, and let $\mathcal R_{\le t}$ be the
lexicographically first basis of $M_t$ for this order. Then:
\begin{enumerate}[label=(\arabic*)]
\item $\mathcal R_{\le1}\subseteq\mathcal R_{\le2}\subseteq\dots\subseteq\mathcal R_{\le T}=B'$;
\item $|\mathcal R_{\le t}|=|P_{\le t}|$, and $E_{\mathcal R_{\le t},P_{\le t}}$ is invertible;
\item every column $j\in P_{\le t}\cup F_{\le t}$ has a nonzero entry in some row of $C\cap\mathcal R_{\le t}$.
\end{enumerate}
\end{lemma}

\begin{proof}
(2). Since $E_P$ is invertible, the columns $P_{\le t}$ are linearly independent, so $\operatorname{rank}M_t=|P_{\le t}|$.
A basis of $M_t$ is therefore exactly a set $\mathcal R$ of $|P_{\le t}|$ rows with $E_{\mathcal R,P_{\le t}}$
invertible.

(1). Suppose the restriction of a row $y$ to $P_{\le t+1}$ lies in the span of the restrictions of a set $X$ of rows.
Then so does its restriction to $P_{\le t}$. Hence $\cl_{M_{t+1}}(X)\subseteq\cl_{M_t}(X)$. A row $y$ lies in the
lexicographically first basis of a matroid $N$ iff $y\notin\cl_N(\text{rows listed before }y)$. Therefore
\[
  y\in\mathcal R_{\le t}\ \Rightarrow\ y\notin\cl_{M_t}(\text{earlier rows})\ \Rightarrow\
  y\notin\cl_{M_{t+1}}(\text{earlier rows})\ \Rightarrow\ y\in\mathcal R_{\le t+1}.
\]
Finally $\mathcal R_{\le T}=B'$, because $M_T$ has rank $r=|B'|$.

(3). Because the rows of $C$ are listed first, the lexicographically first basis contains a maximal independent subset of
$C$. So it maximises $|C\cap\mathcal R|$ among the bases $\mathcal R$ of $M_t$. Let $j\in P_{\le t}\cup F_{\le t}$, and
let $y\in C$ have $E_{y,j}\ne0$; such a $y$ exists because $C$ is a cover. If $y\in\mathcal R_{\le t}$, we are done.
Otherwise
\[
  E_{y,P_{\le t}}=\sum_{w\in\mathcal R_{\le t}}\alpha_wE_{w,P_{\le t}} .
\]
Column $j$ lies in the span of the columns $P_{\le t}$ over all rows. This is trivial for $j\in P_{\le t}$, and holds by
\eqref{eq:A-span} for $j\in F_{\le t}$. Hence $E_{y,j}=\sum_w\alpha_wE_{w,j}$, so some $w\in\mathcal R_{\le t}$ has
$\alpha_wE_{w,j}\ne0$. If $w\in C$, column $j$ is covered. If $w\notin C$, then $\alpha_w\ne0$, so
$\mathcal R_{\le t}-w+y$ is a basis of $M_t$ with one more row of $C$. This contradicts maximality.
\end{proof}

\begin{remark}
Neither minimality of $C$ nor any bound on $|C|$ is used. The order matters: if the rows of $B'\setminus C$ are listed
first, conclusion~(3) can fail. Lemma~\ref{lem:SC} was also confirmed by exact rational rank computations on $43\,947$
random multi-stage configurations with $r\le8$; see Appendix~\ref{app:cert}.
\end{remark}

\subsection{The Staged Volume Lemma}\label{sec:A-svl}

Let $B'$ be a rank-full set of $k$ rows, $C\subseteq B'$ a cover, and $\mathcal R_{\le t}$ the nested bases of
Lemma~\ref{lem:SC}. The \emph{stage rows} are $C_t:=\mathcal R_{\le t}\setminus\mathcal R_{<t}$. For a row $y$ and a
stage $t$ put
\[
  c_y^{(\le t)}:=n_y\Big/\prod_{j\in P_{\le t}}q_j^{E_{yj}},
\]
the product of all primes of $y$ outside $P_{\le t}$, free primes and tail included. The entries of the free columns in
the rows outside $C$ are part of the type $\tau$.

\begin{theorem}[Staged Volume Lemma]\label{thm:svl}\label{thm:SVL}
% v3: Theorem A.5 (SVL)
In the recursion of Lemma~\ref{lem:pil}, integrate the pivots stage by stage, the stage-$t$ pivots $P_t$ by the
change of variables $q_{P_t}\mapsto(n_y)_{y\in C_t}$. Then the following hold.
\begin{enumerate}[label=(\roman*),leftmargin=*]
\item \emph{(Stage identity with the cover split.)} After stage $t$, the Jacobian integrand is
  $\prod_{P_{\le t}}q\big/\bigl(|\det|\prod_{y\in\mathcal R_{\le t}}n_y\bigr)$, and
  \begin{equation}\label{eq:svl-identity}
    \frac{\prod_{j\in P_{\le t}}q_j}{\prod_{y\in\mathcal R_{\le t}}n_y}
    =W_t^{\mathrm{cov}}\prod_{y\in\mathcal R_{\le t}\setminus C}\frac1{n_y}\prod_{j\in P_{\le t}}q_j^{\,1-s_j^{\mathrm{cov}}}.
  \end{equation}
  Here $W_t^{\mathrm{cov}}:=\prod_{y\in C\cap\mathcal R_{\le t}}1/c_y^{(\le t)}$, and
  $s_j^{\mathrm{cov}}:=\sum_{y\in C\cap\mathcal R_{\le t}}E_{yj}\ge1$ by Lemma~\ref{lem:SC}\,(3). Further,
  $|\det|=|\det E_{\mathcal R_{\le t},P_{\le t}}|=\prod_{s\le t}|\det S_s|$, the product of the Schur complements. This
  is a nonzero integer, so $|\det|\ge1$, and the integrand is at most the right side of \eqref{eq:svl-identity}.
  Individual $|\det S_s|$ may be $<1$. The rows outside the cover keep the exact factor $1/n_y$, where $n_y\in J$ is
  their stage integration variable.
\item \emph{(Late primes are harmonic; no type count.)} For $f\in F_t$, Lemma~\ref{lem:SC}\,(3) puts $f$ in a cover row
  of $\mathcal R_{\le t}$, so $W_t^{\mathrm{cov}}$ carries $q_f^{-a}$ with $a\ge1$. The whole column of $f$ is
  determined by its $\mathcal R_{\le t}$-entries:
  \[
    e_f=E_{\cdot,P_{\le t}}\,E_{\mathcal R_{\le t},P_{\le t}}^{-1}\,e_{f,\mathcal R_{\le t}} .
  \]
  Hence the sup-sum of the block is at most $\prod_{y\in C\cap\mathcal R_{\le t}}$ (the $\mathrm{Eu}_s$-ratio of the
  slot range of $F_t$), carried as in Lemma~\ref{lem:fac}\,(S3). The indicators of the remaining rows do not involve $f$.
\item \emph{(Later stages; projected feasibility.)} After the sup-sum of each block, multiply by $\mathbf 1(\mathrm{Proj}_t)$.
  Here $\mathrm{Proj}_t$ is the projection onto the remaining variables of the real completions: summed primes
  $\ge X_0$, and integrated pivots $\ge\tau_j$ in slot order. It is deterministic, and $\mathcal F_{p-}$-measurable at
  the next pivot $p$. On $\mathrm{Proj}_t$, a row $y$ joining at a stage $t'>t$ satisfies
  \begin{equation}\label{eq:B1}
    n_y\ \ge\ c_y^{\mathrm{noF}}\prod_{j\in P_{t'}}q_j^{E_{yj}}\prod_{j\in P_{<t'}}\tau_j^{E_{yj}}\,X_0^{e_y(F_{\le t})}.
  \end{equation}
  Here $c_y^{\mathrm{noF}}$ is $c_y^{(\le t')}$ with the primes of $F_{\le t}$ removed, and $e_y(F_{\le t})$ is the total
  exponent of those primes in $y$. This is the only inequality the later stages use. In particular, the pivot mass of
  later rows carries $\tau_j^{-E}$, not $X_0^{-E}$.
\item \emph{(Slot-indexed chains.)} The Euler chains are indexed by slot, as in Lemma~\ref{lem:fac}\,(S3), so the
  pieces telescope exactly.
\item \emph{(Joint caps.)} A row's free pieces across blocks, together with its tail, take one joint cap:
  $\mathbf 1\bigl(\prod_tc_t\le Y\bigr)\le Y/\prod_tc_t$. Capping block by block would give $Y^{T}$. For rows outside the
  cover the joint cap is not needed: their free part is a monomial in primes of the cover rows, fixed by the type, and
  only their deterministic $X_0$-smooth tail is counted.
\end{enumerate}
\end{theorem}

\begin{proof}
\emph{(i).} We integrate the runs $P_1,P_2,\dots$ in this order, as Lemma~\ref{lem:pil} requires. Suppose stages
$1,\dots,t-1$ are done. The rows of $\mathcal R_{<t}$ have become integration variables $n_y\in J$. The later pivots
$P_{<t}$ are then functions of these variables and of the variables not yet processed; this uses the invertibility of
$E_{\mathcal R_{<t},P_{<t}}$ (Lemma~\ref{lem:SC}\,(2)). In logarithmic coordinates, for $y\in C_t$,
\[
  \log n_y=S_t\,x_{P_t}+(\text{terms in }n_{\mathcal R_{<t}}\text{ and in variables decided before }P_t),
\]
\[
  S_t:=E_{C_t,P_t}-E_{C_t,P_{<t}}E_{\mathcal R_{<t},P_{<t}}^{-1}E_{\mathcal R_{<t},P_t}.
\]
Here $S_t$ is the Schur complement of $E_{\mathcal R_{<t},P_{<t}}$ in $E_{\mathcal R_{\le t},P_{\le t}}$, and it is
invertible because both matrices are. The change of variables $q_{P_t}\mapsto(n_y)_{y\in C_t}$ has Jacobian
$|\det S_t|\prod_{y\in C_t}n_y/\prod_{j\in P_t}q_j$. Multiplying over $s\le t$, and using
$\det E_{\mathcal R_{\le t},P_{\le t}}=\prod_{s\le t}\det S_s$ (block $LU$ factorisation), gives the stated integrand.

For the identity \eqref{eq:svl-identity}, write $n_y=c_y^{(\le t)}\prod_{j\in P_{\le t}}q_j^{E_{yj}}$ for each
$y\in\mathcal R_{\le t}$. For the cover rows, the product is
\[
  \prod_{y\in C\cap\mathcal R_{\le t}}n_y=\prod_{y\in C\cap\mathcal R_{\le t}}c_y^{(\le t)}\cdot\prod_jq_j^{s_j^{\mathrm{cov}}} .
\]
Dividing $\prod_{P_{\le t}}q_j$ by this, and by the factors $n_y$ of the remaining rows of $\mathcal R_{\le t}$, gives
\eqref{eq:svl-identity}. The integer matrix $E_{\mathcal R_{\le t},P_{\le t}}$ is invertible, so
$|\det E_{\mathcal R_{\le t},P_{\le t}}|\ge1$.

\emph{(ii).} By \eqref{eq:A-span}, $e_f=E_{\cdot,P_{\le t}}v_f$ for some $v_f$. Restricting to the rows of
$\mathcal R_{\le t}$ gives $e_{f,\mathcal R_{\le t}}=E_{\mathcal R_{\le t},P_{\le t}}v_f$, hence
$v_f=E_{\mathcal R_{\le t},P_{\le t}}^{-1}e_{f,\mathcal R_{\le t}}$ and the displayed formula. In particular, once the
entries in the rows outside $C$ are fixed by the type, the block data are determined by the cover-row pieces
$(c_{y,F_t})_{y\in C\cap\mathcal R_{\le t}}$, where $c_{y,F_t}$ is the product of the primes of $F_t$ in $y$ with their
exponents. This map is injective.

By Lemma~\ref{lem:SC}\,(3), $f$ has a nonzero entry in a cover row of $\mathcal R_{\le t}$. So
$W_t^{\mathrm{cov}}=\prod_{y\in C\cap\mathcal R_{\le t}}c_{y,F_t}^{-1}\cdot(\text{factors free of }F_t)$, with every
prime of $F_t$ appearing to a power $\ge1$. Before the sup-sum, replace $q_j^{1-s_j^{\mathrm{cov}}}$ by
$\tau_j^{1-s_j^{\mathrm{cov}}}$; this is valid on $\mathrm{Proj}$, since $s_j^{\mathrm{cov}}\ge1$. After that, the only
dependence on $F_t$ is through the cover-row pieces. Hence, by injectivity,
\[
  \sum_{F_t}\prod_{y\in C\cap\mathcal R_{\le t}}c_{y,F_t}^{-1}\le\prod_{y\in C\cap\mathcal R_{\le t}}
  \sum_{c}\frac1c ,
\]
where $c$ runs over the products of realised primes with slots in the range of $F_t$. Each factor is an
$\mathrm{Eu}_s$-ratio, as in Lemma~\ref{lem:fac}\,(S3).

For the remaining rows $y\notin\mathcal R_{\le t}$, the coefficient of $x_f$ in their constraints, after substituting
the pivots of $P_{\le t}$, is
\[
  e_{f,y}-E_{y,P_{\le t}}E_{\mathcal R_{\le t},P_{\le t}}^{-1}e_{f,\mathcal R_{\le t}}=e_{f,y}-E_{y,P_{\le t}}v_f=0 .
\]
So their indicators do not involve $f$. The non-cover rows of $\mathcal R_{\le t}$ carry $1/n_y$ with $n_y$ an
integration variable, which does not involve $f$ either.

\emph{(iii).} After a sup-sum, the summed primes are no longer variables. The remaining integrand is a function of the
variables not yet processed, and every actual configuration lies in $\mathrm{Proj}_t$. So inserting
$\mathbf 1(\mathrm{Proj}_t)$ costs nothing. $\mathrm{Proj}_t$ is defined by real inequalities among deterministic
quantities and variables decided before the next pivot, so it has the stated measurability.

For a row $y$ joining at stage $t'>t$, write
\[
  n_y=c_y^{\mathrm{noF}}\cdot(\text{primes of }F_{\le t}\text{ in }y)\cdot\prod_{j\in P_{t'}}q_j^{E_{yj}}
  \prod_{j\in P_{<t'}}q_j^{E_{yj}} .
\]
The left side, as a function of the remaining variables, does not depend on the completion. On $\mathrm{Proj}_t$ some
completion has all summed primes $\ge X_0$ and all integrated pivots $\ge\tau_j$. This gives \eqref{eq:B1}. In the
stage-$t'$ identity, this inequality replaces the exact factorisation of the rows joining at $t'$; it drops the summed
primes from their cofactors and keeps $\tau_j^{-E_{yj}}$ for their earlier-integrated pivot mass.

\emph{(iv)} is Lemma~\ref{lem:fac}\,(S3). \emph{(v)}: if $\prod_tc_t\le Y$ then $Y/\prod_tc_t\ge1$. For a row outside
the cover, every prime lies in some cover row, and the type fixes its exponents there, so only its tail varies.
\end{proof}

\begin{remark}
As a numerical corroboration only, and not part of the proof, the structural conclusions of Theorem~\ref{thm:svl} were
checked by exact rational arithmetic (computer check; see Appendix~\ref{app:cert}). There was no violation on adversarial
families with $k\le12$ and up to nine stages:
\begin{itemize}
\item $s_j^{\mathrm{cov}}\ge1$;
\item harmonic coverage of late primes;
\item the vanishing of the coefficients of late primes in the remaining rows;
\item the identity \eqref{eq:svl-identity}.
\end{itemize}
\end{remark}

\subsection{Theorem R}\label{sec:A-R}

\begin{theorem}[Rank-full clusters]\label{thm:R}
% v3: Theorem A.6 (Theorem R)
For every window $J\subset[z,2z]$ of length $\delta\le1$ and every $k\ge1$,
\[
  \E\Bigl[\#\{\text{rank-}k\ k\text{-subsets of }\NN\cap J\}\cdot\mathbf 1\bigl(\Good(<t_J)\bigr)\Bigr]
  \le(\theta A_R\delta)^k\,e^{(3H(X_0)+2)k/X_0}\bigl(1+k/X_0\bigr)^{1.02k}\,e^{2.2k/\Lambda},
  \qquad A_R:=1.8(\log z+4).
\]
Since $H(X_0)/X_0\le e^{-10^{7}}$, the factor $e^{(3H(X_0)+2)k/X_0}$ is at most $e^{k\,e^{-10^{6}}}$.
% CORRECTIONS-vs-v3.md entry 1: v3 states e^{2k/X0}.
\end{theorem}

\begin{proof}
\emph{Reduction to patterns.} A $k$-subset of rank $k$ has $k!$ orderings. Fix one ordering. Record:
\begin{itemize}
\item the \emph{pattern}, i.e.\ the matrix $E_P$ of its greedy pivots, with columns in decision order;
\item the slot labels of the pivots;
\item the cofactors $c_1,\dots,c_k$.
\end{itemize}
These determine the subset, since $n_i=c_i\prod_jq_j^{E_{ij}}$. So the map is injective, and the expectation is at most
$1/k!$ times the sum over patterns of the expected number of realisations with that pattern.

\emph{One pattern.} Apply Lemma~\ref{lem:pil}, evaluating the recursion by Theorem~\ref{thm:svl} with the cover $C:=B'$.
Every row is then a cover row. Condition~(3) of Lemma~\ref{lem:SC} holds trivially: pivot columns are nonzero on
$\mathcal R_{\le t}$ by~(2), and late free columns by \eqref{eq:A-span} and~(2). Collect the factors:
\begin{itemize}
\item every row's $n_y$ is integrated over $J$ in exactly one stage, which gives $\delta^k$;
\item $1/|\det E_P|\le1$;
\item the pivot label sums give $\theta^k$;
\item every row carries exactly one harmonic Euler chain over its free pieces, plus one tail sum, and no capped cofactor.
\end{itemize}
For the last item:
\begin{itemize}
\item by Lemma~\ref{lem:fac}\,(S3) and \ref{eq:sF4}, the chain is at most $(\log2z/\Lambda)\,e^{2.20/\Lambda+1/X_0}$;
\item by \ref{eq:sF3p}, the tail sum over $X_0$-smooth $d$ is at most $e^{\gamma}\Lambda(1+1/(2\Lambda^2))$.
\end{itemize}
Their product is at most $A_R$ per row, apart from the chain-end factor $e^{2.2/\Lambda}$ per row, which gives
$e^{2.2k/\Lambda}$.

The remaining factor is $\prod_jq_j^{1-s_j}$. By \eqref{eq:svl-identity} and \eqref{eq:B1}, it is at most
$\prod_j\tau_j^{-(s_j-1)}$, with $s_j$ the full column sums. Summing it over patterns with the factor $1/k!$, by
Lemma~\ref{lem:pattern}, gives $e^{(3H(X_0)+2)k/X_0}(1+k/X_0)^{1.02k}$.
\end{proof}

\begin{remark}
Theorem~\ref{thm:R} contains two simpler cases:
\begin{itemize}
\item tree-like types;
\item types in which every free prime is decided before every pivot. For these, Lemma~\ref{lem:volume} alone gives the
  bound, with the free data summed at the end.
\end{itemize}
\end{remark}

\begin{corollary}[The case of rank at least $r_0$]\label{cor:R-Ecl}
% v3: Theorem A.6, application to E_cl
Let $k=r_0=\lceil L/(3\ell)\rceil+3$ and $\delta=4L^{-6}$, and let $L\ge465$. Then the right side of
Theorem~\ref{thm:R}, multiplied by the number $z/\delta$ of windows at height $z\le x$, is at most $L^{-4}$. This is the
rank-$\ge r_0$ case of Proposition~\ref{prop:ecl}, and it is valid for every $L\ge465$.
\end{corollary}

\begin{proof}
Per element, the factor in Theorem~\ref{thm:R} is
$\theta A_R\delta\,e^{(3H(X_0)+2)/X_0}(1+r_0/X_0)^{1.02}e^{2.2/\Lambda}$. Here:
\begin{itemize}
\item $\theta\,e^{(3H(X_0)+2)/X_0}e^{2.2/\Lambda}\le1.01$, since $H(X_0)/X_0\le e^{-10^{7}}$ and $\theta=1+1/L$;
\item $A_R\le1.8(L+4)$, since $z\le x$;
\item $1+r_0/X_0\le1+L\le2L$.
\end{itemize}
Hence the factor is at most
\[
  1.8(L+4)\cdot4L^{-6}\cdot(2L)^{1.02}\cdot1.01\le15L^{-3.98}\qquad(L\ge465).
\]
Since $z/\delta\le e^{L}L^{6}/4$ and $\log15-3.98\ell<0$,
\[
  \log\Bigl(\frac z\delta\cdot\mathrm{RHS}\Bigr)\le L+6\ell+r_0(\log15-3.98\ell)
  \le L\Bigl(-0.327+\frac{0.91}{\ell}\Bigr)+6\ell+9\le-4\ell
\]
for $L\ge465$. The middle step uses $r_0\ge L/(3\ell)+3$. For the last step, the left side at $L=465$ is $-37.3$, against
$-4\ell=-24.6$, and the difference is increasing in $L$.
\end{proof}
